\documentclass[12pt,a4paper]{article}

\usepackage[margin=1in]{geometry}
\usepackage{amssymb}
\usepackage{booktabs}
\usepackage{hyperref}
\usepackage{enumitem}
\usepackage{titlesec}
\usepackage{fancyhdr}
\usepackage{xcolor}
\usepackage{array}
\usepackage{tabularx}
\usepackage{multirow}
\usepackage{amsmath}
\usepackage{colortbl}
\usepackage{longtable}

% ── Colours ───────────────────────────────────────────────
\definecolor{headblue}{HTML}{1A5276}
\definecolor{highRisk}{HTML}{E74C3C}
\definecolor{medRisk}{HTML}{F39C12}
\definecolor{lowRisk}{HTML}{27AE60}
\definecolor{rowgray}{HTML}{F2F3F4}

% ── Section styling ───────────────────────────────────────
\titleformat{\section}
  {\Large\bfseries\color{headblue}}
  {\thesection.}{0.5em}{}[\vspace{2pt}\titlerule]
\titleformat{\subsection}
  {\large\bfseries\color{headblue!80}}
  {\thesubsection}{0.5em}{}

% ── Header / Footer ──────────────────────────────────────
\pagestyle{fancy}
\setlength{\headheight}{14.5pt}
\addtolength{\topmargin}{-2.5pt}
\fancyhf{}
\lhead{\small\textit{Risk Management --- Traffic Prediction System}}
\rhead{\small\textit{Group E-16}}
\cfoot{\thepage}

\begin{document}

% ═══════════════════════════════════════════════════════════
% TITLE PAGE
% ═══════════════════════════════════════════════════════════
\begin{titlepage}
\centering
\vspace*{2cm}
{\Huge\bfseries\color{headblue} Software Risk Management\\[12pt]}
{\LARGE Risk Identification \& Mitigation Strategies\\[8pt]}
{\large Traffic Prediction and Analysis System\\[20pt]}
\rule{0.6\textwidth}{1pt}\\[20pt]

{\large\textbf{Assignment No. 9}}\\[30pt]

\begin{tabular}{ll}
\textbf{Group:} & E-16 \\[4pt]
\textbf{Members:} & Harshal Jaiswal \\
 & Pranav Jagtap \\
 & Suyash Jobanputra \\
 & Shreyash Jobanputra \\
\end{tabular}

\vfill
{\small Department of Computer Engineering\\
\today}
\end{titlepage}

% ═══════════════════════════════════════════════════════════
\section{Problem Statement}
% ═══════════════════════════════════════════════════════════
Identify potential risks and prepare mitigation strategies for the \textbf{Traffic Prediction and Analysis System} --- a software development project. The exercise involves categorizing risks by type, assessing their probability and impact, constructing a risk matrix, and defining proactive mitigation and contingency plans for each identified risk.

% ═══════════════════════════════════════════════════════════
\section{Introduction}
% ═══════════════════════════════════════════════════════════
Risk management is a systematic process of identifying, analyzing, and responding to uncertainties that may adversely affect a software project's scope, schedule, cost, or quality. In the context of the \textbf{Traffic Prediction and Analysis System} --- a MERN Stack web application that integrates multiple third-party APIs, a MobileNet machine learning model, GPS-based data collection, and a citizen gamification engine --- risk management is especially critical due to the system's reliance on external services and real-time data processing.

The risk management process for this project follows four key steps:
\begin{enumerate}[noitemsep, label=\roman*)]
    \item \textbf{Risk Identification} --- Enumerate all potential risks across technical, project, and external dimensions.
    \item \textbf{Risk Analysis} --- Assess each risk's probability of occurrence and severity of impact.
    \item \textbf{Risk Prioritization} --- Plot risks on a risk matrix to classify them as High, Medium, or Low priority.
    \item \textbf{Risk Mitigation \& Contingency} --- Define proactive strategies to prevent risks and reactive plans to handle them if they occur.
\end{enumerate}

% ═══════════════════════════════════════════════════════════
\section{Risk Identification}
% ═══════════════════════════════════════════════════════════

Risks for the Traffic Prediction and Analysis System are categorized into three groups: Technical, Project/Organizational, and Business/External.

\subsection{Technical Risks}

\begin{center}
\begin{tabular}{c l}
\toprule
\textbf{ID} & \textbf{Risk Description} \\
\midrule
TR-1 & Third-party API failure or rate-limit exceeded (Google Maps, OpenWeatherMap, Calendarific) \\
TR-2 & MobileNet ML model produces inaccurate pothole detection (false positives/negatives) \\
TR-3 & Performance degradation under concurrent prediction requests \\
TR-4 & MongoDB data loss or corruption during development \\
TR-5 & Integration failure between frontend React.js and backend Node.js APIs \\
TR-6 & Cloudinary image upload service outage or cost escalation \\
TR-7 & Browser compatibility issues with Google Maps JavaScript API \\
\bottomrule
\end{tabular}
\end{center}

\subsection{Project / Organizational Risks}

\begin{center}
\begin{tabular}{c l}
\toprule
\textbf{ID} & \textbf{Risk Description} \\
\midrule
PR-1 & Schedule slippage due to underestimated task complexity \\
PR-2 & Reduced team availability during university examination periods \\
PR-3 & Scope creep (uncontrolled addition of features beyond original requirements) \\
PR-4 & Unclear API contract between frontend and backend causing rework \\
PR-5 & Loss of codebase due to lack of version control discipline \\
\bottomrule
\end{tabular}
\end{center}

\subsection{Business / External Risks}

\begin{center}
\begin{tabular}{c l}
\toprule
\textbf{ID} & \textbf{Risk Description} \\
\midrule
BR-1 & Low citizen participation in pothole/complaint reporting (low adoption) \\
BR-2 & Malicious/fraudulent data submissions by unverified citizen reporters \\
BR-3 & Data privacy concerns from GPS-based traffic and location data collection \\
BR-4 & API pricing model changes by Google, OpenWeatherMap, or Calendarific \\
BR-5 & Inaccurate historical traffic data for newly registered routes \\
\bottomrule
\end{tabular}
\end{center}

% ═══════════════════════════════════════════════════════════
\section{Risk Analysis}
% ═══════════════════════════════════════════════════════════

Each risk is evaluated on two dimensions:
\begin{itemize}[noitemsep]
    \item \textbf{Probability (P):} Likelihood of the risk occurring --- rated 1 (Very Low) to 5 (Very High).
    \item \textbf{Impact (I):} Severity of consequences if the risk occurs --- rated 1 (Negligible) to 5 (Catastrophic).
    \item \textbf{Risk Score (RS):} $\text{RS} = P \times I$ --- used to classify the risk level.
\end{itemize}

\bigskip
\begin{center}
\small
\begin{tabular}{c l c c c l}
\toprule
\textbf{ID} & \textbf{Risk} & \textbf{P (1--5)} & \textbf{I (1--5)} & \textbf{RS} & \textbf{Level} \\
\midrule
\rowcolor{highRisk!20}
TR-1 & API Rate Limit / Failure & 4 & 5 & 20 & \textbf{\color{highRisk}High} \\
\rowcolor{medRisk!20}
TR-2 & ML Model Inaccuracy & 3 & 4 & 12 & \textbf{\color{medRisk}Medium} \\
\rowcolor{medRisk!20}
TR-3 & Performance Degradation & 3 & 4 & 12 & \textbf{\color{medRisk}Medium} \\
\rowcolor{medRisk!20}
TR-4 & Database Data Loss & 2 & 5 & 10 & \textbf{\color{medRisk}Medium} \\
\rowcolor{highRisk!20}
TR-5 & API Integration Failure & 4 & 4 & 16 & \textbf{\color{highRisk}High} \\
\rowcolor{lowRisk!20}
TR-6 & Cloudinary Service Outage & 2 & 3 & 6 & \textbf{\color{lowRisk}Low} \\
\rowcolor{lowRisk!20}
TR-7 & Browser Compatibility & 2 & 2 & 4 & \textbf{\color{lowRisk}Low} \\
\rowcolor{highRisk!20}
PR-1 & Schedule Slippage & 4 & 4 & 16 & \textbf{\color{highRisk}High} \\
\rowcolor{highRisk!20}
PR-2 & Team Unavailability (Exams) & 5 & 3 & 15 & \textbf{\color{highRisk}High} \\
\rowcolor{medRisk!20}
PR-3 & Scope Creep & 3 & 4 & 12 & \textbf{\color{medRisk}Medium} \\
\rowcolor{medRisk!20}
PR-4 & Unclear API Contract & 3 & 3 & 9 & \textbf{\color{medRisk}Medium} \\
\rowcolor{lowRisk!20}
PR-5 & Version Control Failure & 2 & 5 & 10 & \textbf{\color{medRisk}Medium} \\
\rowcolor{medRisk!20}
BR-1 & Low Citizen Adoption & 3 & 3 & 9 & \textbf{\color{medRisk}Medium} \\
\rowcolor{highRisk!20}
BR-2 & Fraudulent Data Submissions & 4 & 4 & 16 & \textbf{\color{highRisk}High} \\
\rowcolor{medRisk!20}
BR-3 & Data Privacy Concerns & 3 & 4 & 12 & \textbf{\color{medRisk}Medium} \\
\rowcolor{medRisk!20}
BR-4 & API Pricing Changes & 2 & 4 & 8 & \textbf{\color{medRisk}Medium} \\
\rowcolor{lowRisk!20}
BR-5 & Incomplete Historical Data & 3 & 2 & 6 & \textbf{\color{lowRisk}Low} \\
\bottomrule
\end{tabular}
\end{center}

% ═══════════════════════════════════════════════════════════
\section{Risk Matrix}
% ═══════════════════════════════════════════════════════════

The risk matrix below plots each risk by Probability (rows) and Impact (columns). $\text{RS} = P \times I$: values $\geq 15$ are \textcolor{highRisk}{\textbf{High}}, $8$--$14$ are \textcolor{medRisk}{\textbf{Medium}}, and $\leq 7$ are \textcolor{lowRisk}{\textbf{Low}}.

\bigskip
\begin{center}
\renewcommand{\arraystretch}{1.6}
\begin{tabular}{c | >{\centering\arraybackslash}m{2.8cm} | >{\centering\arraybackslash}m{2.8cm} | >{\centering\arraybackslash}m{2.8cm} | >{\centering\arraybackslash}m{2.8cm}}
 & \textbf{Low Impact (1--2)} & \textbf{Medium Impact (3)} & \textbf{High Impact (4)} & \textbf{Critical Impact (5)} \\
\hline
\textbf{Very High P (5)} & \cellcolor{medRisk!30} & \cellcolor{highRisk!30} PR-2 & \cellcolor{highRisk!30} & \cellcolor{highRisk!30} \\
\hline
\textbf{High P (4)} & \cellcolor{medRisk!30} & \cellcolor{highRisk!30} & \cellcolor{highRisk!30} TR-1, TR-5, BR-2, PR-1 & \cellcolor{highRisk!30} \\
\hline
\textbf{Medium P (3)} & \cellcolor{lowRisk!30} BR-5 & \cellcolor{medRisk!30} BR-1, PR-4 & \cellcolor{medRisk!30} TR-2, TR-3, PR-3, BR-3 & \cellcolor{medRisk!30} \\
\hline
\textbf{Low P (2)} & \cellcolor{lowRisk!30} TR-7 & \cellcolor{lowRisk!30} TR-6 & \cellcolor{medRisk!30} BR-4 & \cellcolor{medRisk!30} TR-4, PR-5 \\
\hline
\textbf{Very Low P (1)} & \cellcolor{lowRisk!30} & \cellcolor{lowRisk!30} & \cellcolor{lowRisk!30} & \cellcolor{lowRisk!30} \\
\hline
\end{tabular}
\end{center}

\bigskip
\noindent\textbf{Legend:}
\textcolor{highRisk}{$\blacksquare$}~\textbf{High Risk} (RS $\geq$ 15)~\quad
\textcolor{medRisk}{$\blacksquare$}~\textbf{Medium Risk} (RS 8--14)~\quad
\textcolor{lowRisk}{$\blacksquare$}~\textbf{Low Risk} (RS $\leq$ 7)

% ═══════════════════════════════════════════════════════════
\section{Risk Mitigation Strategies}
% ═══════════════════════════════════════════════════════════

\subsection{High-Priority Risks (RS $\geq$ 15)}

\subsubsection*{TR-1: API Rate Limit / Failure (RS = 20)}
\textbf{Probability:} High (4) \quad \textbf{Impact:} Critical (5)

\textbf{Proactive Mitigation:}
\begin{itemize}[noitemsep]
    \item Implement server-side response caching (15-minute TTL) for Google Maps and OpenWeatherMap API responses to reduce call frequency.
    \item Use the Outscraper API as a paid fallback for Google Maps data.
    \item Monitor daily API usage against quota limits with automated alerts at 80\% threshold.
\end{itemize}

\textbf{Contingency Plan:}
If an API fails completely, the system defaults to last-cached data for predictions. The affected data factor is excluded from the weighted traffic score calculation and the user is notified via a banner: ``\textit{Live [API Name] data temporarily unavailable. Showing cached data.}''

\bigskip
\subsubsection*{TR-5: API Integration Failure (RS = 16)}
\textbf{Probability:} High (4) \quad \textbf{Impact:} High (4)

\textbf{Proactive Mitigation:}
\begin{itemize}[noitemsep]
    \item Maintain a shared OpenAPI-style API specification document (Swagger/Postman collection) agreed upon by both frontend and backend developers before coding begins.
    \item Implement contract testing using mock API servers during frontend development.
    \item Schedule weekly integration commits to detect mismatches early.
\end{itemize}

\textbf{Contingency Plan:}
Roll back to the last stable API version and conduct a focused sprint to resolve the contract mismatch before resuming feature development.

\bigskip
\subsubsection*{PR-1: Schedule Slippage (RS = 16)}
\textbf{Probability:} High (4) \quad \textbf{Impact:} High (4)

\textbf{Proactive Mitigation:}
\begin{itemize}[noitemsep]
    \item Build a 5-week buffer into the 20-week project schedule beyond the 15-week critical path.
    \item Use aggressive time-boxing: each task has a fixed deadline; incomplete features are de-scoped rather than delayed.
    \item Conduct weekly progress reviews against the Gantt chart milestones.
\end{itemize}

\textbf{Contingency Plan:}
Identify and defer non-critical features (e.g., gamification leaderboard, advanced analytics) to Phase 2. Prioritize core prediction engine and authentication delivery.

\bigskip
\subsubsection*{PR-2: Team Unavailability During Exams (RS = 15)}
\textbf{Probability:} Very High (5) \quad \textbf{Impact:} Medium (3)

\textbf{Proactive Mitigation:}
\begin{itemize}[noitemsep]
    \item Map university exam schedule against the project timeline and schedule low-effort tasks (documentation, testing) during examination periods.
    \item Front-load complex development tasks (T4: Backend, T6: Integration) before the exam window.
    \item Assign flexible non-critical tasks (T3: Database Design) with 2-week float to exam periods.
\end{itemize}

\textbf{Contingency Plan:}
Utilize the 5-week schedule buffer for recovery. The team lead coordinates a catch-up sprint the week after exams conclude.

\bigskip
\subsubsection*{BR-2: Fraudulent Data Submissions (RS = 16)}
\textbf{Probability:} High (4) \quad \textbf{Impact:} High (4)

\textbf{Proactive Mitigation:}
\begin{itemize}[noitemsep]
    \item Implement MobileNet ML validation with a confidence threshold of $\geq 0.85$ for pothole image verification before accepting a report.
    \item Require user authentication (JWT) for all complaint/report submissions.
    \item Apply rate limiting per user account (maximum 5 reports per day).
    \item Implement admin review workflow for flagged submissions before they affect the traffic score.
\end{itemize}

\textbf{Contingency Plan:}
Introduce a temporary manual review queue for all submitted reports if the automated detection confidence drops below acceptable levels. Blacklist accounts with a high rejection rate.

\subsection{Medium-Priority Risks (RS 8--14)}

\begin{center}
\small
\begin{tabularx}{\textwidth}{c X X}
\toprule
\textbf{ID} & \textbf{Mitigation Strategy} & \textbf{Contingency Plan} \\
\midrule
TR-2 & Set MobileNet confidence threshold $\geq 0.85$. Augment training data with diverse pothole images. & Label uncertain detections as ``Pending Review'' and route to admin queue. \\[6pt]
TR-3 & Implement MongoDB indexes on all query fields. Use Node.js clustering for multi-core utilization. Apply 15-min caching for prediction results. & Scale horizontally on Render cloud; reduce prediction frequency during peak load. \\[6pt]
TR-4 & Enable automated daily MongoDB Atlas backups. Use Git for schema migrations. & Restore from the most recent backup. Replay missed transactions from application logs. \\[6pt]
PR-3 & Enforce a formal Change Control Board (CCB) process; any new feature requires written approval from the team lead and project guide. & Defer all unplanned features to a clearly labeled Phase 2 backlog. \\[6pt]
PR-4 & Maintain a Postman API collection as the single source of truth. API contracts frozen before development begins. & 2-day joint debugging sprint between frontend and backend developers to resolve the mismatch. \\[6pt]
BR-1 & Launch with pre-seeded complaint data for demo routes; run a campus awareness campaign to drive initial adoption. & Operate system with admin-curated data until organic citizen participation grows. \\[6pt]
BR-3 & Anonymize GPS coordinates before storage. Publish a clear privacy policy. Do not store personally identifiable raw location data. Add consent prompts. & Halt location-based features and consult legal/academic advisor for guidance. \\[6pt]
BR-4 & Stay within free-tier API limits where possible. Cache aggressively. Monitor cost dashboards weekly. & Switch to lower-cost API alternatives (OpenStreetMap instead of Google Maps if budget exceeded). \\[6pt]
PR-5 & Enforce mandatory Git commits via pull requests on a shared repository. Each member must push daily. Protect the main branch from direct pushes. & Recover from the most recent GitHub remote commit. Use VS Code Timeline for local recovery. \\[6pt]
\bottomrule
\end{tabularx}
\end{center}

\subsection{Low-Priority Risks (RS $\leq$ 7)}

\begin{center}
\small
\begin{tabularx}{\textwidth}{c X X}
\toprule
\textbf{ID} & \textbf{Mitigation Strategy} & \textbf{Contingency Plan} \\
\midrule
TR-6 & Use Cloudinary free tier with local disk as fallback for image storage during development. & Store images on the server filesystem temporarily until Cloudinary is restored. \\[6pt]
TR-7 & Develop and test on Google Chrome (primary). Use CSS autoprefixer to handle cross-browser styling. & Document known limitations. Flag unsupported browsers with a compatibility notice. \\[6pt]
BR-5 & Use the Outscraper API backfill script to supplement missing historical traffic data for new routes. & Substitute city-average traffic score as a fallback when route-specific history is unavailable. \\[6pt]
\bottomrule
\end{tabularx}
\end{center}

% ═══════════════════════════════════════════════════════════
\section{Risk Monitoring Plan}
% ═══════════════════════════════════════════════════════════

Risk monitoring is conducted throughout the project lifecycle (January -- April 2026) using the following mechanisms:

\begin{center}
\begin{tabularx}{\textwidth}{l X l}
\toprule
\textbf{Activity} & \textbf{Description} & \textbf{Frequency} \\
\midrule
Weekly Risk Review & Team reviews the risk register against project progress and Gantt milestones. & Weekly \\[4pt]
API Usage Monitoring & Automated alert at 80\% of daily API quota for Google Maps and OpenWeatherMap. & Daily (automated) \\[4pt]
Integration Smoke Tests & Run a full frontend--backend API smoke test after every major commit to the main branch. & Per major commit \\[4pt]
Attendance Tracking & Track team member availability relative to the project task schedule. & Weekly \\[4pt]
Data Integrity Checks & Run automated scripts to validate MongoDB data consistency and check for anomalous complaint entries. & Twice weekly \\[4pt]
Risk Register Update & Update probability, impact, and status of each risk as the project evolves. & Bi-weekly \\[4pt]
\bottomrule
\end{tabularx}
\end{center}

% ═══════════════════════════════════════════════════════════
\section{Conclusion}
% ═══════════════════════════════════════════════════════════
This assignment identified 17 potential risks across technical, project, and business/external categories for the Traffic Prediction and Analysis System. Risk analysis revealed five high-priority risks (RS $\geq$ 15): API failure (TR-1, RS=20), integration failure (TR-5, RS=16), schedule slippage (PR-1, RS=16), fraudulent data submissions (BR-2, RS=16), and team unavailability during exams (PR-2, RS=15).

Proactive mitigation strategies --- including response caching, shared API specifications, time-boxed scheduling, ML-based fraud detection, and mandatory version control practices --- address the root causes of these risks. Contingency plans ensure continuity of development even if a risk materializes.

By maintaining a living risk register and conducting weekly reviews against the project timeline, the development team can respond to emerging risks in a structured and timely manner, ensuring the successful delivery of the Traffic Prediction and Analysis System within the planned 15-week development window.

\end{document}
